\documentclass[conference]{IEEEtran}
\usepackage{booktabs,amsmath,url,array}
\usepackage[hidelinks]{hyperref}
\title{TRACE-RAG: Chronological Trust, Corroboration and Retroactive Remediation for Knowledge-Base Poisoning}
\author{Anonymous preprint draft --- author information to be completed}
\date{November 2026}
\begin{document}
\maketitle
\begin{abstract}
Retrieval-augmented generation (RAG) systems trust retrieved text, so a small
amount of poisoned content can alter an answer long before a static corpus
audit notices. We describe TRACE-RAG, an integrated defence that treats
ingestion as a chronological stream: it scores passages with cheap signals,
maintains document-, family- and source-level trust, corroborates suspicious
claims against independent evidence, defers unresolved release, and re-opens
answers whose cited content is quarantined later. We implement twelve attack
families and eleven policies behind one evaluator contract, with explicit
answer-level attack-success accounting and missing-data markers. This draft
reports an \emph{offline fixture} matrix, not a benchmark result. It completed
132/132 cases, but detection precision and recall were zero and the run used a
hashing embedder, stub generator and lexical verifier. The real-model Natural
Questions study is specified and instrumented but has not been executed. We
therefore present the design, protocol and negative/plumbing evidence, and make
no state-of-the-art or deployment claim.
\end{abstract}
\begin{IEEEkeywords}
retrieval-augmented generation, knowledge-base poisoning, trust, provenance,
corroboration, retroactive remediation, reproducibility
\end{IEEEkeywords}
\section{Introduction}
RAG grounds generation in retrieved passages, but retrieval is an untrusted
input channel. PoisonedRAG demonstrated that a few adversarial passages can
steer answers \cite{poisonedrag}. Defences have since explored trust filtering,
certified aggregation, forensic tracing and provenance. The open engineering
question is not whether a detector can be built, but whether trust can be
maintained over time: evidence arrives in a stream, a passage may look benign
when released, and an answer may need correction after later evidence arrives.
TRACE-RAG addresses that lifecycle. Its contributions are (i) a chronological
evaluation contract shared by twelve attacks and eleven policies; (ii) explicit
answer-level ASR that cannot be inferred from exposure; (iii) integrated
suspicion, trust, corroboration, deferred release and retroactive remediation;
and (iv) a reproducible protocol separating fixture evidence from a real-model
study. The current evidence is deliberately reported as incomplete.
\section{Related Work}
PoisonedRAG establishes knowledge-corruption attacks \cite{poisonedrag}.
TrustRAG filters retrieval using trust signals \cite{trustrag}; RobustRAG
aggregates isolated responses with certificates \cite{robustrag}. RAGForensics
traces poisoning sources \cite{ragforensics}; FilterRAG filters retrieval
\cite{filterrag}; RAGDefender detects and mitigates poisoned knowledge
\cite{ragdefender}. SafeRAG and CorruptRAG provide security benchmarks and
data \cite{saferag,corruptrag}. RAGShield is a recent provenance-oriented
defence \cite{ragshield}. Our closest distinction is the combination of
provenance-independent corroboration, deferred release and retroactive answer
remediation with blast-radius accounting. We do not claim novelty over every
individual component, and the adapted baselines in this draft are not the
published algorithms.
\section{Threat Model}
An attacker can insert passages through one or more sources at chosen stream
steps. It may know the retrieval pipeline (white-box), adapt after quarantine,
split corroboration across Sybil sources, or coordinate near-duplicate claims.
The defender observes ingestion order and provenance, but not evaluator labels.
Target queries and evaluator-only gold/attacker answers are fixed before the
run. The defender may abstain; abstention is counted as attack failure, never as
attack success. Out of scope: model-weight poisoning, prompt injection in the
user turn, and attacks that never enter the corpus.
\section{System}
The pipeline indexes each chunk with provenance, computes six cheap suspicion
signals, and stores a feature snapshot. A logistic scorer with isotonic
calibration maps signals to a suspicion score. Trust is tracked per document,
near-duplicate family and source. HIGH-suspicion or refuted content is blocked
or quarantined; unresolved medium-suspicion content is withheld until
corroboration. Trust-weighted retrieval uses
$\mathrm{score}=\mathrm{sim}\cdot\max(t_{\mathrm{eff}},\epsilon)^\lambda$.
Every answer records its context and citations. When later evidence quarantines
a cited document, the ledger identifies prior answers that relied on it and
emits remediation records. Labels are used only by the offline trainer, never
by the live policy.
\section{Experimental Protocol}
\subsection{Attacks and policies}
The twelve attacks are entity swap, negation, instruction injection, framing,
signal evasion, hit-and-run, slow burn, black box, white box, Sybil, coordinated
and adaptive. All are deterministic given a seed and emit the same ingestion
event as benign content. Policies are no defence, duplicate filter, n-gram
perplexity surrogate, suspicion-only, corroboration-only, trust threshold, trust
without deferred verification, adapted TrustRAG, adapted RobustRAG, always-on
leave-one-out verification, and full RAGLedger. Adapted methods are marked
\texttt{published\_method=False}.
\subsection{Data and execution tiers}
Tier 1 is the bundled mini corpus with a hashing embedder, stub generator and
lexical verifier. Tier 2 is a 5,000-passage Natural Questions pilot on a Colab
GPU. Tier 3 is a multi-seed NQ/BEIR benchmark with Contriever/FAISS and an NLI
verifier. Only Tier 1 has been executed. Tiers 2--3 are specified in
\texttt{docs/RESEARCH\_PROTOCOL.md} and the Colab notebook; their absence is a
limitation, not an implied result.
\subsection{Metrics and statistics}
Primary outcomes are targeted poison exposure and answer-level ASR on eligible
queries. Detection precision/recall/F1 and clean-passage false-positive rate are
reported separately. p99 latency and peak traced memory are secondary. Each
field has a measured/not-measured status. The planned real study uses three to
five independent seeds, paired policy comparisons, query-clustered bootstrap
intervals and Holm correction. A single-seed fixture cannot support inferential
claims.
\section{Results}
Table~\ref{tab:offline} reports the measured fixture. All 132 cases completed
without execution failures. Mean exposure over attacks fell from 0.833 without
defence to 0.208 with RAGLedger; mean answer ASR fell from 0.444 to 0.139.
RAGLedger reached zero exposure for adaptive, black-box, coordinated and
signal-evasion attacks, but retained 0.333 exposure for seven families. These
are plumbing observations, not benchmark efficacy.
egin{table}[t]
\centering
\caption{Offline fixture matrix. Values are means over one seed and three targeted queries per case; the scope is 	exttt{SMOKE TEST} with a hashing embedder, stub generator, and lexical verifier. These are plumbing results, not benchmark evidence.}
\label{tab:offline}
egin{tabular}{lrrrrrrrr}
	oprule
Attack & \multicolumn{4}{c}{No defence} & \multicolumn{4}{c}{RAGLedger} \\
\cmidrule(lr){2-5}\cmidrule(lr){6-9}
 & Exp. & ASR & P & R & Exp. & ASR & P & R \\
\midrule
adaptive & 0.667 & 0.000 & 0.000 & 0.000 & 0.000 & 0.000 & 0.000 & 0.000 \\
black\_box & 1.000 & 1.000 & 0.000 & 0.000 & 0.000 & 0.000 & 0.000 & 0.000 \\
coordinated & 1.000 & 1.000 & 0.000 & 0.000 & 0.000 & 0.000 & 0.000 & 0.000 \\
entity\_swap & 1.000 & 0.667 & 0.000 & 0.000 & 0.333 & 0.333 & 0.000 & 0.000 \\
framing & 1.000 & 0.000 & 0.000 & 0.000 & 0.333 & 0.000 & 0.000 & 0.000 \\
hit\_and\_run & 1.000 & 0.667 & 0.000 & 0.000 & 0.333 & 0.333 & 0.000 & 0.000 \\
instruction\_injection & 0.667 & 0.000 & 0.000 & 0.000 & 0.167 & 0.000 & 0.000 & 0.000 \\
negation & 1.000 & 0.000 & 0.000 & 0.000 & 0.333 & 0.000 & 0.000 & 0.000 \\
signal\_evasion & 0.667 & 0.000 & 0.000 & 0.000 & 0.000 & 0.000 & 0.000 & 0.000 \\
slow\_burn & 1.000 & 0.667 & 0.000 & 0.000 & 0.333 & 0.333 & 0.000 & 0.000 \\
sybil & 1.000 & 0.667 & 0.000 & 0.000 & 0.333 & 0.333 & 0.000 & 0.000 \\
white\_box & 1.000 & 0.667 & 0.000 & 0.000 & 0.333 & 0.333 & 0.000 & 0.000 \\
ottomrule
\end{tabular}
\end{table}

The important negative result is that detection precision and recall were zero
for every group. The HIGH-suspicion band produced no labeled true positives in
this fixture. Exposure reduction therefore cannot be attributed to validated
detection; it may reflect policy exclusion or retrieval effects. RQ2 is not
supported. No claim about trust-signal quality is warranted.
\section{Discussion and Limitations}
The prototype demonstrates that twelve attacks, eleven policies, answer-level
ASR, memory/latency and artifact export can run through one chronological
contract. It does not demonstrate real-model robustness, utility preservation,
statistical significance, or superiority over published defences. The fixture
uses a stub generator and lexical verifier; its numbers must not be compared
with NQ or BEIR results. The adapted TrustRAG/RobustRAG policies are
simplifications, and the detector's zero recall is a concrete blocker.
The real study must additionally report clean-answer withholding, abstention,
retroactive remediation precision and blast radius. If detection remains at
zero recall, the trust stage must be revised before any efficacy claim.
\section{Reproducibility}
The repository pins the notebook to an implementation commit, installs with
\texttt{pip install -e ".[all,data]"}, verifies a real CUDA matmul and asserts
that \texttt{trace\_rag} is imported from the checkout. Fixture results are in
\texttt{baseline\_logs/research\_smoke\_full/}; the full matrix command is in
\texttt{docs/RESEARCH\_PROTOCOL.md}. The notebook's heavy flags default to
false and fail visibly when prerequisites are absent. Generated \texttt{runs/}
artifacts are not committed; the analysis JSON/CSV and logs are.
\section{Conclusion}
TRACE-RAG provides an integrated, chronological defence and a reproducible
evaluation harness. The current evidence supports only the plumbing claim: the
matrix runs end to end, but detection failed in the fixture and real-model
experiments remain unexecuted. We release the protocol and negative evidence so
the next run can test the hypothesis rather than inherit an unverified claim.
\begin{thebibliography}{99}
\bibitem{poisonedrag} X.~Zou \emph{et al.}, ``PoisonedRAG: Knowledge Corruption Attacks to Retrieval-Augmented Generation of Large Language Models,'' arXiv:2402.07867, 2024. \url{https://arxiv.org/abs/2402.07867}
\bibitem{trustrag} ``TrustRAG: Enhancing Robustness and Trustworthiness in RAG,'' arXiv:2501.00879, 2025. \url{https://arxiv.org/abs/2501.00879}
\bibitem{robustrag} ``RobustRAG: Certified Generation in Retrieval-Augmented Generation,'' arXiv:2405.15556, 2024. \url{https://arxiv.org/abs/2405.15556}
\bibitem{ragforensics} ``RAGForensics: Tracing Poisoning Sources in Retrieval-Augmented Generation,'' arXiv:2504.21668, 2025. \url{https://arxiv.org/abs/2504.21668}
\bibitem{filterrag} ``FilterRAG: Robust Retrieval-Augmented Generation against Knowledge Poisoning,'' arXiv:2508.02835, 2025. \url{https://arxiv.org/abs/2508.02835}
\bibitem{ragdefender} ``RAGDefender: Detecting and Mitigating Knowledge Poisoning in RAG,'' arXiv:2511.01268, 2025. \url{https://arxiv.org/abs/2511.01268}
\bibitem{saferag} ``SafeRAG: Benchmarking Security in Retrieval-Augmented Generation,'' arXiv:2501.18636, 2025. \url{https://arxiv.org/abs/2501.18636}
\bibitem{corruptrag} ``CorruptRAG: A Large-Scale Dataset for Security Evaluation of RAG,'' arXiv:2504.03957, 2025. \url{https://arxiv.org/abs/2504.03957}
\bibitem{ragshield} ``RAGShield: A Provenance-Based Defence for Retrieval-Augmented Generation,'' arXiv:2604.00387, 2026. \url{https://arxiv.org/abs/2604.00387}
\end{thebibliography}
\end{document}
